\subsection{Quasi-projective varieties}

\bdefn

    A {\emphcolor graded ring} is a ring $S$ together with a decomposition $S=\bigoplus_{i\geq0}S_i$, where each $S_i\subseteq S$ is an Abelian group and $S_i\cdot S_j\subseteq S_{i+j}$.
    In particular $S_0\subseteq S$ is a subring, and $S_i\subseteq S$ is an $S_0$-submodule, and the direct sum $S_+=\bigoplus_{i>0}S_i$ is an ideal of $S$.

    Elements of $S_i$ are called {\emphcolor homogeneous of degree $i$}, and an element is just {\emphcolor homogeneous} if it is homogeneous of some degree.
    Denote the set of homogeneous elements of $S$ by $S_\homo$.

    An ideal $I\leq S$ is {\emphcolor homogeneous} if it respects the decomposition: $I=\bigoplus_{i\geq0}(I\cap S_i)$.

\edefn

\bexam

    Let $S=k[x_1,\dots,x_n]$, then $S$ is graded by $S_d=\lspan_k\set{x_1^{i_1}\cdots x_n^{i_n}}[i_1+\cdots+i_n=d]$, called the space of homogeneous polynomials of degree $d$.

\eexam

For $S$ graded and $a\in S$, we can uniquely write $a$ as $a=\sum_da_d$ for $a_d\in S_d$, all but finitely many being zero.
The element $a_d$ is called the {\emphcolor homogeneous part of $a$ of degree $d$}.
Clearly $(a+b)_d=a_d+b_d$ so the map $S\to S_d$ mapping an element to its homogeneous part of degree $d$ is an Abelian group morphism.
Furthermore,
$$ \parens{\sum_da_d}\parens{\sum_db_d} = \sum_d\sum_ka_kb_{d-k} , \qquad \sum_ka_kb_{d-k}\in S_d, $$
so that $(ab)_d=\sum_ka_kb_{d-k}$.

We let $\deg(a)$ denote the maximal $d$ such that $a_d\neq0$.
So we see that if $\deg(a)\leq n$ and $\deg(b)\leq m$, then $\deg(ab)\leq n+m$ and $(ab)_{n+m}=a_nb_m$.

And clearly an ideal $I$ is homogeneous iff for all $a\in I$ we have $a_d\in I$.

\blemm

    \benum
        \item An ideal $I\leq S$ is homogeneous iff it is generated by homogeneous elements, so $I=(I\cap S_\homo)$.
        \item Homogeneous ideals are closed under sums, products, intersections, and radicals.
        \item A non-unit homogeneous ideal $I\leq S$ is prime iff for all $a,b\in S_\homo$ such that $a,b\notin I$ we have $ab\notin I$.
        \item A finitely generated homogeneous ideal is generated by finitely many homogeneous elements.
    \eenum

\elemm

\bproof

    \benum
        \item If $I\leq S$ is homogeneous, then
        $$ I = \sum_i(I\cap S_i) = \parens{\bigcup I\cap S_i} = (I\cap S_\homo) . $$
        Conversely suppose $(b_i)$ generate $I$, where $b_i\in S_{d_i}$.
        Then elements of $I$ are of the form $a=\sum_ia_ib_i$ with $a_i\in S$.
        Then
        $$ a_d = \sum_i(a_ib_i)_d = \sum_i(a_i)_{d-d_i}b_i \in I , $$
        so $I$ is homogeneous.

        \item This is clear from (1) for sums and products.
        And if $a\in I=\bigcap_iI_i$ then $a_d\in I_i$ for all $i$ so $a_d\in I$ and thus $I$ is homogeneous.

        Now if $I$ is homogeneous, we induct on $d$ to show that for every $f\in\sqrt I$ such that $f=\sum_{i=0}^df_i$, $f_i\in\sqrt I$.
        For $d=0$ we have $f=f_0$ so this is clear.
        Now inductively, it is sufficient to show that $f_d\in\sqrt I$ since then $f-f_d=\sum_{i=0}^{d-1}f_i\in\sqrt I$ and we proceed by induction.

        By definition there is some $n$ such that $f^n\in I$.
        Now, we know $\deg(f)\leq d$ and thus inductively $\deg(f^n)\leq dn$ so $(f^n)_{dn}=f_d^n$.
        Since $f^n\in I$ we have $f_d^n=(f^n)_{dn}\in I$ and thus $f_d\in\sqrt I$ as desired.

        \item One direction is clear.
        Now suppose we have that $ab\notin I$ when $a,b\in S_\homo-I$.
        For $a,b\in S-I$ we have $a=\sum_ia_i$ and $b=\sum_jb_j$.
        Let us write $a_I=\sum_{i,a_i\in I}a_i$ and $a^I=\sum_{i,a_i\notin I}a_i$ so $a=a_I+a^I$ and similarly $b=b_I+b^I$.
        So $a_I,b_I\in I$ and
        $$ ab = a_Ib_I + a_Ib^I + a^Ib_I + a^Ib^I , $$
        so it is sufficient to show that $a^Ib^I\notin I$.

        That is, it is sufficient to consider the case when $a_i\notin I$ and $b_i\notin I$ to show $ab\notin I$.
        Let $n=\deg(a)$ and $m=\deg(b)$ so $(ab)_{n+m}=a_nb_m$, and since $a_n,b_m\in S_\homo-I$ we have $a_nb_m\notin I$, and by homogeneity $ab\notin I$.

        \item Let $T\subseteq I$ be finite and generate $I$.
        Then let
        $$ T_\homo = \set{a_i}[a\in T,i\geq0] . $$
        Since $I$ is homogeneous, $T_\homo\subseteq I$.
        And since elements of $T$ are finite sums of elements from $T_\homo$, we have $I=(T)\subseteq(T_\homo)\subseteq I$ so $(T_\homo)=I$.
        Clearly $T_\homo$ is a finite subset of $S_\homo$.
        \qed
    \eenum

\eproof

For an algebraically closed field $k$, the group $k^\times$ acts on $\bA^n=k^n$ via multiplication: for $P=(a_1,\dots,a_n)\in k^n$ and $\lambda\in k^\times$ we have the action $\lambda P=(\lambda a_1,\dots,\lambda a_n)$.

For the graded ring $S=k[\bA^n]$ with grading $S_d=\lspan_k\set{x^I}[\sum I=d]$, we see that for $f\in S_d$ we have $f(\lambda P)=\lambda^df(P)$.
This is clear for monomials $x^I$, and follows in general by linearity.

\blemm

    A closed subset $Z\subseteq\bA^n$ is invariant under $k^\times$'s action iff $\I(Z)\subseteq S$ is homogeneous.
    For a general subset $X\subseteq\bA^n$, if $X$ is $k^\times$-invariant then $\I(X)\leq S$ is homogeneous (only one direction of the previous equivalence holds).

\elemm

\bproof

    Suppose $\I(Z)\subseteq S$ is homogeneous, and let $P\in Z$ and $\lambda\in k^\times$.
    Since $Z$ is closed, $Z=Z(\I(Z))$ it is sufficient to show that for $f\in\I(Z)$ we have $f(\lambda P)=0$.
    If $f$ is homogeneous, this is clear as $f(\lambda P)=\lambda^df(P)=0$.
    Then in general we have $f=\sum_df_d$ and $f_d\in\I(Z)$ so $f_d(\lambda P)=0$ and thus $f(\lambda P)=0$.

    Now if $X$ is $k^\times$-invariant, we want to show that if $f\in\I(X)$ then $f_d\in\I(X)$.
    Fix $P\in X$, so we want to show $f_d(P)=0$ for all $d$.
    For every $\lambda\in k^\times$, by invariance we have $\lambda P\in X$ so $f(\lambda P)=0$.
    Thus
    $$ 0 = f(\lambda P) = \sum_df_d(\lambda P) = \sum_d\lambda^df_d(P) . $$
    That is, the polynomial $Q(t)=\sum_df_d(P)x^d$ satisfies $Q(\lambda)=0$ for all $\lambda\in k^\times$.
    So $Q$ has infinitely many roots, and thus must be the zero polynomial, so $f_d(P)=0$ for all $d$, as desired.
    \qed

\eproof

Notice that since $\cl X=Z(\I(X))$ and thus $\I(\cl X)=\I(X)$, we see that if $X$ is $k^\times$-invariant so too is $\cl X$.

We view $S=k[\bA^{n+1}]$ as the graded $k$-algebra $k[x_0,\dots,x_n]$, so we begin indexing at $0$ here.
Since $\bA^{n+1}-0$ is $k^\times$-invariant, we can quotient by $k^\times$ to get $\bP^n_k=\bP^n$.
That is, we define $\bP^n$ to be the set of all $k^\times$-orbits of $\bA^{n+1}-0$.
We denote the orbit of $(a_0,\dots,a_n)\in\bA^{n+1}-0$ by $[a_0:\dots:a_n]\in\bP^n$.
Note that equivalently, $\bP^n$ can be viewed as the set of all lines in $\bA^{n+1}$ passing through the origin, and then $[a_0:\dots:a_n]$ corresponds to the line $t(a_0,\dots,a_n)$.

Let $\pi\colon\bA^{n+1}-0\to\bP^n$ be the canonical projection map: $\pi(a_0,\dots,a_n)=[a_0:\dots:a_n]$.
This is surjective, and thus we can endow $\bP^n$ with the quotient topology.
In particular this makes $\bP^n$ into an SWF $(\bP^n,\O_\reg)$ where
$$ \O_\reg(U) = \set{f\colon U\to k}[f\circ\pi\in\O_\reg(\pi^{-1}U)] . $$
That is, $f\colon U\to k$ is regular iff the pullback $\pi^*(f)\colon\pi^{-1}U\to k$ is regular.

\bdefn

    For a homogeneous polynomial $f\in S_\homo$ and $P\in\bP^n$ we write $f(P)=0$ to mean that $f(\tilde P)=0$ for some representative (equivalently every representative) $\tilde P\in\pi^{-1}(P)$.
    For $f\in S_\homo$ let
    $$ Z_{\bP^n}(f) = \set{P\in\bP^n}[f(P)=0] , $$
    and in general for $T\subseteq S_\homo$ let
    $$ Z_{\bP^n}(T) = \bigcap_{f\in T}Z_{\bP^n}(f) . $$
    For an homogeneous ideal $I\leq S$ we define $Z_{\bP^n}(I)=Z_{\bP^n}(I\cap S_\homo)$.
    If $T$ is any of the above (an homogeneous polynomial/set/ideal) then we define $D_{\bP^n}(T)=\bP^n-Z_{\bP^n}(T)$.

    Finally for any $Y\subseteq\bP^n$ we define the homogeneous ideal $\I_\homo(Y)\leq S$ by
    $$ \I_\homo(Y) = \gen{f\in S_\homo}[\hbox{$f(P)=0$ for all $P\in Y$}] . $$
    As it is generated by homogeneous elements, it is an homogeneous ideal.

\edefn

\blemm

    \benum
        \item $Z_{\bP^n}$ is order reversing,
        \item $Z_{\bP^n}(0)=\bP^n$ and $Z_{\bP^n}(1)=\emptyset$,
        \item $Z_{\bP^n}(T_1\cdot T_2)=Z_{\bP^n}(T_1)\cup Z_{\bP^n}(T_2)$ and $Z_{\bP^n}\parens{\bigcup_iT_i}=\bigcap_iZ_{\bP^n}(T_i)$,
        \item For every $T\subseteq S_\homo$, there is an equality $Z_{\bP^n}(T)=Z_{\bP^n}((T))$.
    \eenum

\elemm

\bproof

    Points (1) and (2) are clear.
    Point (3) is shown in essentially the same way as in the affine case.
    For point (4) we have
    $$ Z_{\bP^n}((T)) = Z_{\bP^n}((T)\cap S_\homo) , $$
    and $T\subseteq(T)\cap S_\homo$ so we have an inclusion $Z_{\bP^n}((T))\subseteq Z_{\bP^n}(T)$.
    Conversely let $P\in Z_{\bP^n}(T)$ and $f\in(T)\cap S_\homo$ then $f=\sum_ia_it_i$ for $a_i\in S$ and $t_i\in T$.
    Suppose $\deg(f)=d$ then $\deg(a_it_i)=d$ for all $i$ which means that each $a_i$ must be homogeneous as well and $\deg(a_i)=d-\deg(t_i)$.
    Since everything is homogeneous, since $t_i(P)=0$ for all $i$ we have $f(P)=0$ so we have the other inclusion as well.
    \qed

\eproof

\bprop

    \benum
        \item For $T\subseteq S_\homo$ we have the equality
        $$ \pi^{-1}Z_{\bP^n}(T) = Z_{\bA^{n+1}-0}(T) . $$
        \item $Z\subseteq\bP^n$ is closed iff $Z=Z_{\bP^n}(T)$ for some $T\subseteq S_\homo$.
        \item The collection $\set{D_{\bP^n}(f)}_{f\in S_\homo}$ forms a basis of open sets of $\bP^n$.
    \eenum

\eprop

\bproof

    We know that for $f\in S_\homo$ and $\tilde P\in\bA^{n+1}-0$, $f(\tilde P)=0$ iff $f(\pi(\tilde P))=0$.
    Thus $\pi(\tilde P)\in Z_{\bP^n}(T)$ iff $f(\pi(\tilde P))=0$ for all $f\in T$ iff $f(\tilde P)=0$ for all $f\in T$ iff $\tilde P\in Z_{\bA^{n+1}-0}(T)$ as desired.

    The preimage $\pi^{-1}Z_{\bP^n}(T)=Z_{\bA^{n_1}-0}(T)$ is closed in $\bA^{n+1}-0$.
    Since $\bP^n$ has the quotient topology, this means $Z_{\bP^n}(T)$ is closed in $\bP^n$.
    Conversely if $Z\subseteq\bP^n$ is closed, then $\pi^{-1}Z\subseteq\bA^{n+1}-0$ is closed and $k^\times$-invariant.
    Since $\cl{\pi^{-1}Z}=Z_{\bA^{n+1}}(\I(\pi^{-1}Z))$, we see that $I=\I(\pi^{-1}Z)$ is homogeneous.
    So $I=(I\cap S_\homo)$ and thus $\cl{\pi^{-1}Z}=Z_{\bA^n}(I)=Z_{\bA^n}(I\cap S_\homo)$ so we have equality
    $$ \pi^{-1}Z = \cl{\pi^{-1}Z}\cap(\bA^{n+1}-0) = Z_{\bA^{n+1}-0}(I\cap S_\homo) = \pi^{-1}Z_{\bP^n}(I\cap S_\homo) . $$
    Thus we see $Z=Z_{\bP^n}(I\cap S_\homo)$ as desired.
    The reverse direction is clear since $\pi^{-1}Z_{\bP^n}(T)=Z_{\bA^{n+1}-0}(T)$ is closed.

    Every closed set in $\bP^n$ is of the form $Z_{\bP^n}(T)=\bigcap_{f\in T}Z_{\bP^n}(f)$ so all open sets are unions of $D_{\bP^n}(f)$.
    \qed

\eproof

For homogeneous polynomials of the same degree $g,h\in S_d$ and $\tilde P\in D_{\bA^{n+1}}(h)$ and $\lambda\in k^\times$ we see
$$ \frac gh(\lambda\tilde P) = \frac{g(\lambda\tilde P)}{h(\lambda\tilde P)} = \frac{\lambda^dg(\tilde P)}{\lambda^dh(\tilde P)} = \frac{g(\tilde P)}{h(\tilde P)} . $$
So homogeneous rational functions $g/h$ define maps $D_{\bP^n}(h)\to k$.

This allows us to better describe the SWF $(\bP^n,\O_\reg)$.

\blemm

    For $U\subseteq\bP^n$ open, a function $f\colon U\to k$ is regular iff there are homogeneous polynomials $g,h\in S_d$ such that $U\subseteq D_{\bP^n}(h)$ and $f=g/h|_U$.

\elemm

\bproof

    We know that $f$ is regular iff $\pi^*(f)\colon\pi^{-1}U\to k$ is regular.
    This is iff it is rational, so there are polynomials $g,h\in S$ such that $\pi^{-1}U\subseteq D(h)$ and for everu $\tilde P\in\pi^{-1}U$ we have
    $$ \pi^*(f)(P) = f(\pi(P)) = \frac{g(\tilde P)}{h(\tilde P)} . $$
    In particular, for $g,h\in S_d$ the function $f=g/h$ is regular.

    Conversely we want to show that two polynomials satisfying the above equality must be homogeneous of the same degree.
    First we claim that $\pi^{-1}U\subseteq D(h)$ implies $h$ homogeneity.
    Suppose not, so there are $i_1\neq i_2$ such that $h_{i_1},h_{i_2}$ are nonzero.
    We have open non-empty subsets $D(h_{i_1}),D(h_{i_2}),\pi^{-1}U\subseteq\bA^{n+1}$, and since $\bA^{n+1}$ is irreducible this means that their intersection is nonempty.
    So there is $\tilde P\in\pi^{-1}U$ such that $h_{i_1}(\tilde P),h_{i_2}(\tilde P)\neq0$.

    Then for every $\lambda\in k^\times$ we have $\lambda\tilde P\in\pi^{-1}U\subseteq D(h)$ so $h(\lambda\tilde P)\neq0$.
    Hence
    $$ 0\neq h(\lambda\tilde P) = \sum_ih_i(\lambda\tilde P) = \sum_i\lambda^ih_i(\tilde P) , $$
    then as before this means $Q(t)=\sum_ih_i(\tilde P)t^i$ has no zeroes in $k^\times$.
    Since $k$ is algebraically closed, this means its only zero is $0$ and thus $Q(t)=at^d$ for some $a\in k$.
    This means that $h_i(\tilde P)=0$ for all $i\neq d$, so there is at most a single index where $h_i(\tilde P)\neq0$, a contradiction.

    Now we claim that $g$ is homogeneous.
    For every $\tilde P\in\pi^{-1}U$ and $\lambda\in k^\times$ we have
    $$ \frac{g(\lambda\tilde P)}{\lambda^dh(\tilde P)} = \frac{g(\lambda\tilde P)}{h(\lambda\tilde P)} = f(\pi(\lambda\tilde P)) = f(\pi(\tilde P)) = \frac{g(\tilde P)}{h(\tilde P)} . $$
    Thus we see $g(\lambda\tilde P)=\lambda^dg(\tilde P)$.
    Thus for all $\lambda\in k^\times$,
    $$ \sum_i\lambda^ig_i(\tilde P) = \lambda^d\sum_ig_i(\tilde P) . $$
    Defining $Q(t)=t^d\sum_ig_i(\tilde P)-\sum_it^ig_i(\tilde P)$ we see that this has infinitely many zeroes and is thus identically zero, so $g_i(\tilde P)=0$ for all $i\neq d$.
    So $g(\tilde P)=g_d(\tilde P)$ for all $\tilde P\in\pi^{-1}U$.
    Since $\bA^{n+1}$ is Noetherian, $\pi^{-1}U$ is dense and so $g=g_d$ on all $\bA^{n+1}$ as desired.
    \qed

\eproof

\bnote

    The above results, namely that $Z\subseteq\bP^n$ is closed iff it is of the form $Z_{\bP^n}(T)$ and $f\colon U\to k$ is regular iff it is homogeneous rational, are many times taken to be definitions
    of projective space.
    We had to work a little harder to show that the more obvious definition, induced from affine space, agrees with this.

\enote

\bcoro

    Let $Y\subseteq\bP^n$ be an arbitrary subspace, so $(\bP^n,\O_\reg)$ induces an SWF $(Y,\O_\reg)$.
    \benum
        \item The closed sets of $Y$ are those of the form $Z_Y(T)=\set{P\in Y}[f(P)=0\hbox{ for all $f\in T$}]$.
        \item The collection $D_Y(f)=\set{P\in Y}[f(P)\neq0]$ forms a basis of open sets of $Y$.
        \item A function $f\colon Y\to k$ is regular iff it is locally given by homogeneous rational functions.
        That is, for every $P\in Y$ there is an open subset $P\in U\subseteq\bP^n$ such that $f|_U$ is homogeneous rational.
    \eenum

\ecoro

\bproof

    Points (1) and (2) are clear from the definition of the subspace topology, the obvious equalities
    $$ Z_Y(T) = Z_{\bP^n}(T)\cap Y,\qquad D_Y(T) = D_{\bP^n}(T)\cap Y , $$
    and the results we proved above.
    Point (3) is immediate from the definition of induced SWFs and the previous result.
    \qed

\eproof

\blemm

    The SWF $(\bP^n,\O_\reg)$ is an algebraic variety.

\elemm

\bproof

    Cover $\bA^{n+1}-0$ by $\set{D(x_i)}_{i=0}^n$, so through $\pi$ we cover $\bP^n$ by $\set{D_{\bP^n}(x_i)}_{i=0}^n$.
    And clearly we have
    $$ U_i = D_{\bP^n}(x_i) = \set{[a_0:\dots:a_n]\in P}[a_i\neq0] . $$
    We claim that these $U_i$ are affine.

    Define $\phi\colon U_i\to\bA^n$ by
    $$ \phi[a_0:\dots:a_n] = \parens{\frac{a_0}{a_i},\dots,\hat{\frac{a_i}{a_i}},\dots,\frac{a_n}{a_i}} . $$
    Each map $[a_0:\dots:a_n]\mapsto a_j/a_i$ is regular, as it is homogeneous rational (given by $x_j/x_i$), so this is a morphism into $\bA^n$.
    $\phi$ is bijective with inverse
    $$ \psi(a_1,\dots,a_n) = \bracks{a_1:\dots:a_i:1:a_{i+1}:\dots:a_n} . $$
    To show that $\psi$ is a morphism of SWFs, we note that it is equal to the composition
    $$ \bA^n \xtos{\tilde\psi} \bA^{n+1} \xtos\pi \bP^n , $$
    where $\tilde\psi$ maps $(a_1,\dots,a_n)$ to $(a_1,\dots,a_i,1,a_{i+1},\dots,a_n)$, and this is clearly a morphism of SWFs as its components are regular (polynomial even).
    Thus $\psi$ is the composition of morphisms of SWFs and thus a morphism itself.
    \qed

\eproof

\bdefn

    A {\emphcolor projective variety} is an SWF $X$ which is isomorphic to an SWF $(Y,\O_\reg)$ where $Y\subseteq\bP^n$ is closed.

    A {\emphcolor quasi-projective variety} is an SWF $X$ which is isomorphic to an open subspace of an open subvariety of a projective variety.

\edefn

Clearly projective varieties are quasi-projective.
Notice that quasi-projective varieties are just SWFs which are isomorphic to locally closed subvarieties of $\bP^n$.

\bcoro

    \benum
        \item Quasi-projective varieties are algebraic varieties.
        \item Every locally closed subvariety of a quasi-projective variety is quasi-projective.
        \item Every quasi-affine variety is quasi-projective.
    \eenum

\ecoro

\bproof

    The first point follows from the fact that locally closed subvarieties of an algebraic variety are algebraic.
    A locally closed subvariety of a quasi-projective variety is isomorphic to a locally closed subvariety of a locally closed subvariety of $\bP^n$.
    Local closure is a transitive property, so this is just a locally closed subvariety of $\bP^n$ and thus quasi-projective, showing (2).

    Since $\bA^n\cong U_0\subseteq\bP^n$, we see that affine space $\bA^n$ is projective.
    Quasi-affine varieties are locally closed subvarieties of $\bA^n$, and thus locally closed subvarieties of a projective variety, and thus quasi-affine by (2).
    \qed

\eproof

